%!TEX program = xelatex




\documentclass[
  11pt,
  twoside,
  openright
]{book}


\usepackage{polyglossia}
\setdefaultlanguage{english}

\usepackage{fontspec}
\setmainfont{TeX Gyre Termes}[
  BoldFont   = TeX Gyre Termes Bold,
  ItalicFont = TeX Gyre Termes Italic
]
\setsansfont{TeX Gyre Heros}[
  BoldFont   = TeX Gyre Heros Bold,
  ItalicFont = TeX Gyre Heros Italic
]

\setmonofont{TeX Gyre Cursor}[
  Scale    = 0.88,
  BoldFont = TeX Gyre Cursor Bold
]

\usepackage{geometry}
\geometry{
  papersize  = {7in, 10in},
  top        = 1in,
  bottom     = 1.1in,
  inner      = 1in,
  outer      = 0.875in,
  headheight = 14pt,
  twoside
}

\usepackage{xcolor}
\definecolor{BookBlue}       {HTML}{1A4A8A}
\definecolor{BookDark}       {HTML}{1C1C1E}
\definecolor{BookMid}        {HTML}{4A4A4A}
\definecolor{LearnBg}        {HTML}{EEF7EE}
\definecolor{LearnBorder}    {HTML}{27AE60}
\definecolor{TipBg}          {HTML}{EAF4FB}
\definecolor{TipBorder}      {HTML}{2980B9}
\definecolor{WarningBg}      {HTML}{FEF9E7}
\definecolor{WarningBorder}  {HTML}{F39C12}
\definecolor{NoteBg}         {HTML}{F4F4F4}
\definecolor{NoteBorder}     {HTML}{888888}
\definecolor{TakeawayBg}     {HTML}{F0F0FF}
\definecolor{TakeawayBorder} {HTML}{5B4FBE}
\definecolor{CodeBg}         {HTML}{F7F7F7}
\definecolor{CodeComment}    {HTML}{6A737D}
\definecolor{CodeKeyword}    {HTML}{D73A49}
\definecolor{CodeString}     {HTML}{032F62}
\definecolor{CodeConstant}   {HTML}{005CC5}
\definecolor{CodeAnnotation} {HTML}{6A737D}

\usepackage{microtype}
\usepackage{setspace}
\setstretch{1.18}

\usepackage{csquotes}
\usepackage{parskip}
\setlength{\parskip}{6pt}
\setlength{\parindent}{0pt}

\widowpenalty=10000
\clubpenalty=10000
\displaywidowpenalty=10000
\predisplaypenalty=10000
\raggedbottom
\emergencystretch=1.5em
\hbadness=3000
\vbadness=3000

\usepackage{needspace}
\usepackage{etoolbox}
\preto\section{\needspace{5\baselineskip}}
\preto\subsection{\needspace{4\baselineskip}}
\preto\subsubsection{\needspace{3\baselineskip}}

\usepackage{titlesec}
\renewcommand{\bottomtitlespace}{0.15\textheight}

\titleformat{\part}[display]
  {\centering\sffamily\Huge\bfseries\color{BookBlue}}
  {\sffamily\large\mdseries\partname\ \thepart}{20pt}{}

\titleformat{\chapter}[display]
  {\sffamily\huge\bfseries\color{BookBlue}}
  {\sffamily\normalsize\mdseries\chaptertitlename\ \thechapter}{12pt}{}
\titlespacing*{\chapter}{0pt}{0pt}{24pt}

\titleformat{\section}
  {\sffamily\Large\bfseries\color{BookBlue}}
  {\thesection}{1em}{}
\titlespacing*{\section}{0pt}{18pt}{6pt}

\titleformat{\subsection}
  {\sffamily\large\bfseries\color{BookDark}}
  {\thesubsection}{1em}{}
\titlespacing*{\subsection}{0pt}{14pt}{4pt}

\titleformat{\subsubsection}
  {\sffamily\normalsize\bfseries\color{BookDark}}
  {\thesubsubsection}{1em}{}
\titlespacing*{\subsubsection}{0pt}{10pt}{2pt}

\usepackage{fancyhdr}
\pagestyle{fancy}
\fancyhf{}
\fancyhead[LE]{\small\sffamily\color{BookMid}\leftmark}
\fancyhead[RO]{\small\sffamily\color{BookMid}\rightmark}
\fancyfoot[LE,RO]{\small\sffamily\color{BookMid}\thepage}
\renewcommand{\headrulewidth}{0.3pt}
\renewcommand{\footrulewidth}{0pt}

\fancypagestyle{plain}{
  \fancyhf{}
  \fancyfoot[LE,RO]{\small\sffamily\color{BookMid}\thepage}
  \renewcommand{\headrulewidth}{0pt}
}

\makeatletter
\renewcommand{\cleardoublepage}{%
  \clearpage
  \if@twoside
    \ifodd\c@page
    \else
      \thispagestyle{empty}%
      \hbox{}%
      \newpage
      \if@twocolumn\hbox{}\newpage\fi
    \fi
  \fi
}
\makeatother

\usepackage{rotating}
\usepackage{booktabs}
\usepackage{tabularx}
\usepackage{array}
\usepackage{longtable}
\renewcommand{\arraystretch}{1.35}

\usepackage{graphicx}
\usepackage{float}
\usepackage{caption}
\captionsetup{
  font      = {small,sf},
  labelfont = {bf,color=BookBlue},
  skip      = 6pt
}

\usepackage{listings}

\lstset{
  backgroundcolor   = \color{CodeBg},
  basicstyle        = \ttfamily\fontsize{9.5}{11}\selectfont,
  keywordstyle      = \color{CodeKeyword}\bfseries,
  stringstyle       = \color{CodeString}\itshape,
  commentstyle      = \color{CodeComment}\itshape,
  numbers           = left,
  numberstyle       = \tiny\color{BookMid},
  numbersep         = 6pt,
  breaklines        = true,
  frame             = single,
  framerule         = 0.4pt,
  rulecolor         = \color{NoteBorder},
  showstringspaces  = false,
  tabsize           = 2,
  captionpos        = b,
  breakatwhitespace = false,
  columns           = fullflexible,
  linewidth         = \linewidth,
  postbreak         = {},
  xleftmargin       = 1.5em,
  framexleftmargin  = 1.5em,
  aboveskip         = 6pt,
  belowskip         = 4pt,
}

\lstdefinelanguage{rego}{
  keywords={package, import, default, allow, deny, not, as, if, else, some, every, in, contains},
  keywords=[2]{true, false, null},
  keywordstyle=[2]\color{gray}\bfseries,
  sensitive=true,
  morecomment=[l]{\#},
  morestring=[b]",
  morestring=[b]',
  commentstyle=\color{gray}\itshape,
  stringstyle=\color{orange},
  alsoletter={-},
  literate={:=}{{$\gets$}}1 {<-}{{$\leftarrow$}}1
}

\lstdefinelanguage{json}{
  morestring=[b]",
  morestring=[b]',
  morecomment=[l]//,
  morecomment=[s]{/*}{*/},
  literate=
    *{0}{{\color{CodeString}0}}{1}
     {1}{{\color{CodeString}1}}{1}
     {2}{{\color{CodeString}2}}{1}
     {3}{{\color{CodeString}3}}{1}
     {4}{{\color{CodeString}4}}{1}
     {5}{{\color{CodeString}5}}{1}
     {6}{{\color{CodeString}6}}{1}
     {7}{{\color{CodeString}7}}{1}
     {8}{{\color{CodeString}8}}{1}
     {9}{{\color{CodeString}9}}{1}
     {:}{{\color{CodeKeyword}{:}}}{1}
     {,}{{\color{CodeKeyword}{,}}}{1}
     {\{}}{{\color{CodeKeyword}{\{}}} {1}
     {\}}}{{\color{CodeKeyword}{\}}}} {1}
     {[}{{\color{CodeKeyword}{[}}}{1}
     {]}{{\color{CodeKeyword}{]}}}{1},
  sensitive=true,
  alsoletter={:,-},
  morekeywords={true,false,null},
  keywordstyle=\color{CodeKeyword}\bfseries,
  stringstyle=\color{CodeString},
  commentstyle=\color{CodeComment}\itshape,
}

\lstdefinelanguage{TypeScript}{
  morekeywords={abstract,any,as,asserts,async,await,bigint,boolean,
    break,case,catch,class,const,continue,declare,default,delete,do,
    else,enum,export,extends,false,finally,for,from,function,get,
    global,if,implements,import,in,infer,instanceof,interface,is,
    keyof,let,module,namespace,new,null,number,object,of,package,
    private,protected,public,readonly,require,return,satisfies,set,
    static,string,super,switch,symbol,throw,true,try,type,typeof,
    undefined,unique,unknown,using,var,void,while,with,yield,
    override,accessor,out},
  morecomment=[l]{//},
  morecomment=[s]{/*}{*/},
  morestring=[b]",
  morestring=[b]',
  sensitive=true,
  alsoletter={.},
  keywordstyle=\color{CodeKeyword}\bfseries,
  stringstyle=\color{CodeString},
  commentstyle=\color{CodeComment}\itshape,
  breaklines=true,
  showstringspaces=false,
}

\lstdefinelanguage{CSharp}{
  morekeywords={abstract,async,await,base,bool,break,byte,case,catch,
    char,checked,class,const,continue,decimal,default,delegate,do,
    double,else,enum,event,explicit,extern,false,finally,fixed,float,
    for,foreach,goto,if,implicit,in,int,interface,internal,is,lock,
    long,namespace,new,null,object,operator,out,override,params,
    private,protected,public,readonly,ref,return,sbyte,sealed,short,
    sizeof,stackalloc,static,string,struct,switch,this,throw,true,
    try,typeof,uint,ulong,unchecked,unsafe,ushort,using,virtual,void,
    volatile,while,add,alias,ascending,descending,dynamic,from,get,
    global,group,into,join,let,orderby,partial,remove,select,set,
    value,var,where,yield},
  morecomment=[l]{//},
  morecomment=[s]{/*}{*/},
  morestring=[b]",
  morestring=[b]',
  sensitive=true,
  alsoletter={.},
  keywordstyle=\color{CodeKeyword}\bfseries,
  stringstyle=\color{CodeString},
  commentstyle=\color{CodeComment}\itshape,
  breaklines=true,
  breakatwhitespace=false,
  breakindent=0pt,
  columns=fullflexible,
  keepspaces=true,
  postbreak=\mbox{\textcolor{gray}{$\hookrightarrow$}\space},
  showstringspaces=false,
}

\lstdefinelanguage{Java}{
  morekeywords={
    abstract,assert,boolean,break,byte,case,catch,char,class,const,continue,
    default,do,double,else,enum,extends,final,finally,float,for,goto,if,
    implements,import,instanceof,int,interface,long,native,new,package,
    private,protected,public,return,short,static,strictfp,super,switch,
    synchronized,this,throw,throws,transient,try,void,volatile,while,
    module,requires,exports,opens,uses,provides,to,with,transitive,
    var,
    record,sealed,permits,non-sealed,yield
  },
  morekeywords=[2]{true,false,null},
  morekeywords=[3]{@interface},
  sensitive=true,
  morecomment=[l]{//},
  morecomment=[s]{/*}{*/},
  morestring=[b]",
  morestring=[b]',
  alsoletter={.},
  keywordstyle=\color{CodeKeyword}\bfseries,
  keywordstyle=[2]\color{CodeConstant},
  keywordstyle=[3]\color{CodeAnnotation},
  stringstyle=\color{CodeString},
  commentstyle=\color{CodeComment}\itshape,
  breaklines=true,
  breakatwhitespace=false,
  breakindent=0pt,
  columns=fullflexible,
  keepspaces=true,
  postbreak=\mbox{\textcolor{gray}{$\hookrightarrow$}\space},
  showstringspaces=false,
}

\lstdefinelanguage{dockerfile}{
  morecomment=[l]{\#},
  morestring=[b]",
  morestring=[b]',
  morekeywords={FROM,MAINTAINER,RUN,CMD,LABEL,EXPOSE,ENV,ADD,COPY,
    ENTRYPOINT,VOLUME,USER,WORKDIR,ARG,ONBUILD,STOPSIGNAL,HEALTHCHECK,
    SHELL},
  sensitive=true,
  keywordstyle=\color{CodeKeyword}\bfseries,
  commentstyle=\color{CodeComment}\itshape,
  stringstyle=\color{CodeString},
  breaklines=true,
  showstringspaces=false,
}

\usepackage{amsmath}
\lstdefinelanguage{yaml}{
  morecomment=[l]{\#},
  morestring=[b]",
  morestring=[b]',
  moredelim=[l][{\color{CodeKeyword}\bfseries}]{:},
  moredelim=*[l][{\color{CodeString}}]{- },
  moredelim=[s][{\color{CodeKeyword}\bfseries}]{\{}{\}},
  moredelim=[s][{\color{CodeKeyword}\bfseries}]{[}{]},
  literate=
    *{true}{{\color{CodeKeyword}true}}{4}
    {false}{{\color{CodeKeyword}false}}{5}
    {null}{{\color{CodeKeyword}null}}{4},
  sensitive=true,
  keywordstyle=\color{CodeKeyword}\bfseries,
  stringstyle=\color{CodeString},
  commentstyle=\color{CodeComment}\itshape,
  breaklines=true,
  breakatwhitespace=true,
  columns=fullflexible,
  showstringspaces=false
}

\DeclareRobustCommand{\code}[1]{\texttt{#1}}
\providecommand{\term}[1]{\textit{#1}}
\providecommand{\chapterquote}[2]{%
  \vspace{1em}
  \begin{quote}\itshape#1\end{quote}
  \begin{flushright}--- #2\end{flushright}
  \vspace{1em}
}

\usepackage{tcolorbox}
\tcbuselibrary{skins,breakable}

\newtcolorbox{learningbox}{
  colback=LearnBg, colframe=LearnBorder,
  before skip = 12pt, fonttitle=\sffamily\bfseries\small, title={Learning Objective},
  breakable, lines before break=3,
  left=8pt, right=8pt, top=6pt, bottom=6pt,
  boxrule=0.6pt, width=\linewidth
}
\newtcolorbox{tipbox}{
  colback=TipBg, colframe=TipBorder,
  before skip = 12pt, fonttitle=\sffamily\bfseries\small, title={Tip},
  breakable, lines before break=3,
  left=8pt, right=8pt, top=6pt, bottom=6pt,
  boxrule=0.6pt, width=\linewidth
}
\newtcolorbox{warningbox}{
  colback=WarningBg, colframe=WarningBorder,
  before skip = 12pt, fonttitle=\sffamily\bfseries\small, title={Warning},
  breakable, lines before break=3,
  left=8pt, right=8pt, top=6pt, bottom=6pt,
  boxrule=0.6pt, width=\linewidth
}
\newtcolorbox{notebox}{
  colback=NoteBg, colframe=NoteBorder,
  before skip = 12pt, fonttitle=\sffamily\bfseries\small, title={Note},
  breakable, lines before break=3,
  left=8pt, right=8pt, top=6pt, bottom=6pt,
  boxrule=0.6pt, width=\linewidth
}
\newtcolorbox{takeawaybox}{
  colback=TakeawayBg, colframe=TakeawayBorder,
  before skip = 12pt, fonttitle=\sffamily\bfseries\small, title={Key Takeaway},
  breakable, lines before break=3,
  left=8pt, right=8pt, top=6pt, bottom=6pt,
  boxrule=0.6pt, width=\linewidth
}

\usepackage{amssymb}

\PassOptionsToPackage{hyphens}{url}
\usepackage{hyperref}
\usepackage{imakeidx}
\makeindex[columns=2, title=Index, intoc=true]
\urlstyle{tt}
\makeatletter
\g@addto@macro\UrlBreaks{\do\.\do\_\do\-}
\makeatother
\DeclareUrlCommand{\code}{\urlstyle{tt}}

\hypersetup{
  colorlinks  = true,
  linkcolor   = BookBlue,
  urlcolor    = BookBlue,
  citecolor   = BookBlue,
  pdftitle    = {Engineering AI Agents for Production},
  pdfauthor   = {Haythem Aber},
  pdfsubject  = {Engineering AI Agents for Production - From LLMs to Agent Ecosystems, Volume~2},
  pdfkeywords = {AI agents production, Java agent framework, Kubernetes agents, observability evaluation, LLM security governance, production case studies}
}

\gdef\theBookTitle{Engineering AI Agents for Production}
\gdef\theBookSubtitle{Infrastructure, Observability, Security,
and Production Case Studies}
\gdef\theBookAuthor{Haythem Aber}
\gdef\theBookYear{2026}
\gdef\theBookDate{April 2026}
\gdef\theBookVolume{Volume 2}
\gdef\theBookVolumeScope{From LLMs to Agent Ecosystems}

\newcommand{\BookTitle}[1]{\gdef\theBookTitle{#1}}
\newcommand{\BookSubtitle}[1]{\gdef\theBookSubtitle{#1}}
\newcommand{\BookAuthor}[1]{\gdef\theBookAuthor{#1}}
\newcommand{\BookYear}[1]{\gdef\theBookYear{#1}}
\newcommand{\BookDate}[1]{\gdef\theBookDate{#1}}
\newcommand{\BookVolume}[1]{\gdef\theBookVolume{#1}}
\newcommand{\BookVolumeScope}[1]{\gdef\theBookVolumeScope{#1}}

\newcommand{\makecover}{%
  \begin{titlepage}%
    \thispagestyle{empty}%
    \vspace*{2cm}%
    \begin{center}%
      {\sffamily\large\scshape\theBookVolume\par}%
      \vspace{0.3cm}%
      {\sffamily\small\itshape\theBookVolumeScope\par}%
      \vspace{1.2cm}%
      \hrule height 0.5pt\par%
      \vspace{1.2cm}%
      {\sffamily\fontsize{38}{46}\bfseries\selectfont
        \theBookTitle\par}%
      \vspace{1.2cm}%
      \hrule height 0.5pt\par%
      \vspace{0.8cm}%
      {\sffamily\fontsize{11}{14}\selectfont
        \theBookSubtitle\par}%
      \vspace{1.5cm}%
      {\sffamily\Large\theBookAuthor\par}%
      \vfill%
      {\sffamily\normalsize\theBookDate\par}%
      \vspace{0.5cm}%
      {\sffamily\small\itshape
        Volume~1: \textit{Architecting AI Agent Systems}
        --- Available Now\par}%
      \vspace{0.8cm}%
    \end{center}%
  \end{titlepage}%
  \nopagecolor%
}

\newcommand{\partclose}[1]{%
  \cleardoublepage
  \vspace*{\fill}
  \begin{center}
    \sffamily\Large\color{BookBlue}#1
  \end{center}
  \vspace*{\fill}
  \cleardoublepage
}

\newcommand{\makecopyright}{%
  \thispagestyle{empty}%
  \vspace*{\fill}%
  \begin{small}%
    \noindent\textit{\theBookTitle}\\%
    \textit{\theBookSubtitle}\\%
    \textit{\theBookVolume: \theBookVolumeScope}\\[12pt]%
    \noindent Copyright \textcopyright{} \theBookYear\
    \theBookAuthor. All rights reserved.\\[8pt]%
    \noindent No part of this publication may be reproduced,
    distributed, or transmitted in any form or by any means
    without prior written permission of the author.\\[8pt]%
    \noindent\textbf{Permissions:} For rights and licensing
    enquiries, contact the author directly.\\[8pt]%
    \noindent First edition. \theBookDate.%
  \end{small}%
  \vspace*{\fill}%
  \clearpage%
}

\newcommand{\makededication}[1]{%
  \cleardoublepage%
  \thispagestyle{empty}%
  \vspace*{\fill}%
  \begin{center}\itshape #1\end{center}%
  \vspace*{\fill}%
  \clearpage%
}

\newcommand{\makehalftitle}{%
  \cleardoublepage%
  \thispagestyle{empty}%
  \vspace*{\fill}%
  \begin{center}%
    {\sffamily\LARGE\bfseries\color{BookBlue}\theBookTitle}%
  \end{center}%
  \vspace*{\fill}%
  \clearpage%
}

\newenvironment{preface}{%
  \cleardoublepage%
  \chapter*{Preface to Volume~2}%
  \addcontentsline{toc}{chapter}{Preface to Volume~2}%
  \markboth{Preface to Volume~2}{Preface to Volume~2}%
}{\clearpage}

\pretocmd{\tableofcontents}{\cleardoublepage}{}{}
\setcounter{tocdepth}{1}

% ============================================================
%%% === MACROS ===
% ============================================================

\newcommand{\component}[1]{\textsf{\textbf{#1}}}
\newcommand{\seechapter}[1]{\textit{(See Chapter~\ref{#1})}}
\newcommand{\seepart}[1]{\textit{(See Part~\ref{#1})}}

\usepackage{tikz}
\usetikzlibrary{arrows.meta}

\usepackage{pdflscape}
\usepackage{ragged2e}

% ============================================================
\begin{document}
% ============================================================

\BookTitle{Engineering AI Agents for Production}
\BookSubtitle{Infrastructure, Observability, Security,
and Production Case Studies}
\BookVolume{Volume 2}
\BookVolumeScope{From LLMs to Agent Ecosystems}
\BookAuthor{Haythem Aber}
\BookYear{2026}
\BookDate{April 2026}

% ============================================================
\frontmatter
% ============================================================

\makehalftitle
\makecover
\makecopyright
\makededication{%
  Architecture is the set of decisions\\
  that are hard to reverse.\\
  Make them deliberately.%
}

\begin{preface}
% TODO: preface text (preserved from original)
\end{preface}

\tableofcontents

% ============================================================
\mainmatter
% ============================================================

% ============================================================
% PART 1 — Building an Agent Framework from Scratch
% (~130 pages, Chapters 1–4)
% ============================================================
\part{Building an Agent Framework from Scratch}
\label{part:framework}

% ----------------------------------------------------------
% Chapter 1 — DONE: written, reviewed, all fixes applied.
\input{Tex/chapter01}

% ----------------------------------------------------------
\chapter{Architecting the Core Runtime}
\label{ch:core-runtime}
% TODO: Rewrite per edit plan Section 4, Chapter 2.
% Keep: state transition table (Table 2.1), StepExecutionException design, interrupt narrative, two-phase model.
% Delete: ~480 lines of complete class listings (Run, DefaultRunScheduler, RunHandleImpl, InMemoryRepository, OpenAiModelClient, AgentBuilder).
% Replace: each listing → focused excerpt ≤40 lines (transition(), checkpoint(), LEGAL_TRANSITIONS; switch dispatch only).
% Add: "What the Loop Does Not Do" sidebar + Fig 2.1 (sequence) + Fig 2.2 (interrupt state diagram).

% ----------------------------------------------------------
\chapter{Implementing Memory, Tools, and Orchestration}
\label{ch:memory-tools-orchestration}
% TODO: Rewrite per edit plan Section 4, Chapter 3.
% Keep: MemoryRecord contract (~12l), TokenWindowMemoryPolicy explanation, structured-concurrency paragraph,
%        ReactOutputParser.parse() (~25l), McpToolAdapter concept.
% Delete: ~620 lines (3 backend classes, DefaultToolRegistry, ToolExecutionExecutor, McpClient, ReActOrchestrator, AgentBuilder).
% Replace: backends → comparison table; registry → 20-line register()+execute() excerpt; ReAct → 30-line dispatch.
% Add: "Memory Tiers: What the Run Sees vs. What Persists" subsection + Fig 3.1 + Fig 3.2 + Fig 3.3.
% Apply: H3 fix (TaintActionValidator BLOCK behaviour); A1 fix (memory injection pipeline in §9.5).

% ----------------------------------------------------------
\chapter{Implementing the Protocol Stack}
\label{ch:protocol-stack}
% TODO: Rewrite per edit plan Section 4, Chapter 4.
% Keep: TenantPolicyEngine.evaluate() 30-line excerpt, StructuredEventExporter route method 25l, integration tests (expanded).
% Delete: ~580 lines (HttpAgentAdapter 120l, McpServerAdapter 200l, A2AServerAdapter 180l, A2AClient 80l).
% Replace: each adapter → 15-line excerpt showing only the AgentRuntime.run() bridge.
% Add: "Why Protocol Separation Matters" opening + Fig 4.1 (protocol adapter pattern).

\partclose{End of Part~1 --- The framework is built.
Parts~2 through~4 make it production-ready.}

% ============================================================
% PART 2 — Infrastructure for Distributed Agent Systems
% (~140 pages, Chapters 5–11)
% ============================================================
\part{Infrastructure for Distributed Agent Systems}
\label{part:infrastructure}

% ----------------------------------------------------------
\chapter{The Agent Control Plane}
\label{ch:control-plane}
% TODO: Rewrite per edit plan Section 5, Chapter 5.
% Keep: desired-state/observed-state explanation, 3 mandatory health dimensions, 4 workflow patterns taxonomy.
% Delete: Java impl subsections (~60 lines) → replace with config YAML.
% Add: "What the Part 1 Framework Cannot Do Alone" opening + Fig 5.1 (control plane object model).

% ----------------------------------------------------------
\chapter{Sandboxing and Isolation}
\label{ch:sandboxing}
% TODO: Rewrite per edit plan Section 5 standard rule.
% Delete: all listings >30 lines.
% Keep: comparison tables, decision frameworks, architectural diagrams.
% Replace: listings → config/YAML excerpts.
% Add: "What the Part 1 Framework Cannot Do Alone" opening + one architectural diagram.

% ----------------------------------------------------------
\chapter{Messaging and Event-Driven Coordination}
\label{ch:messaging}
% TODO: Rewrite per edit plan Section 5 standard rule.
% Delete: all listings >30 lines.
% Keep: comparison tables, decision frameworks, architectural diagrams.
% Replace: listings → Kafka topic configs, Temporal workflow definitions.
% Add: "What the Part 1 Framework Cannot Do Alone" opening + Fig 8.1 (event-driven workflow patterns).

% ----------------------------------------------------------
\chapter{Event Sourcing for Agent State}
\label{ch:event-sourcing}
% TODO: Rewrite per edit plan Section 5 standard rule.
% Delete: all listings >30 lines.
% Keep: comparison tables, decision frameworks.
% Replace: listings → config excerpts.
% Add: "What the Part 1 Framework Cannot Do Alone" opening + Fig 9.1 (event sourcing sequence diagram).

% ----------------------------------------------------------
\chapter{Deploying Agents on Kubernetes}
\label{ch:kubernetes}
% TODO: Rewrite per edit plan Section 5 standard rule.
% Delete: all listings >30 lines.
% Keep: comparison tables, decision frameworks, architectural diagrams.
% Replace: listings → Kubernetes YAML manifests.
% Add: "What the Part 1 Framework Cannot Do Alone" opening + one architectural diagram.

% ----------------------------------------------------------
\chapter{Multi-Cloud and Hybrid Deployments}
\label{ch:multi-cloud}
% TODO: Rewrite per edit plan Section 5 standard rule.
% Delete: all listings >30 lines.
% Keep: comparison tables, decision frameworks, architectural diagrams.
% Replace: listings → config excerpts.
% Add: "What the Part 1 Framework Cannot Do Alone" opening + one architectural diagram.

% ----------------------------------------------------------
\chapter{Sustainability and Cost Efficiency}
\label{ch:sustainability}
% TODO: Rewrite per edit plan Section 5 standard rule.
% Delete: all listings >30 lines.
% Keep: comparison tables, decision frameworks.
% Replace: listings → config excerpts.
% Add: "What the Part 1 Framework Cannot Do Alone" opening + one architectural diagram.

\partclose{End of Part~2 --- The framework scales.
Part~3 makes it observable, measurable, and testable.}

% ============================================================
% PART 3 — Observability, Evaluation, and Testing
% (~110 pages, Chapters 12–17)
% ============================================================
\part{Observability, Evaluation, and Testing}
\label{part:observability}

% ----------------------------------------------------------
\chapter{Observability: Tracing, Logging, and Metrics}
\label{ch:observability}
% TODO: Rewrite per edit plan Section 6, Chapter 12.
% Keep: 3-level observability hierarchy explanation, metric naming conventions table.
% Delete: duplicate StructuredEventExporter listing (already in Ch 4).
% Add: "Diagnosing a Silent Loop" walkthrough + Fig 12.1 (observability signal hierarchy).

% ----------------------------------------------------------
\chapter{Evaluation Frameworks and Quality Gates}
\label{ch:evaluation}
% TODO: Rewrite per edit plan Section 6 standard rule (Ch 13).
% Delete: listings. Keep: tables, decision frameworks. Add: one diagram.

% ----------------------------------------------------------
\chapter{Explainability and Trajectory Reconstruction}
\label{ch:explainability}
% TODO: Rewrite per edit plan Section 6, Chapter 14 — EXPANDED treatment.
% Show CycleRecord + AgentEvent streams → ExplainabilityFormatter → human-readable trajectory.
% Add: Fig 14.1 (reasoning chain reconstruction flowchart).

% ----------------------------------------------------------
\chapter{Simulation and Synthetic Testing}
\label{ch:simulation}
% TODO: Rewrite per edit plan Section 6 standard rule (Ch 15).
% Delete: listings. Keep: tables, decision frameworks. Add: one diagram.

% ----------------------------------------------------------
\chapter{Replay and Deterministic Debugging}
\label{ch:replay}
% TODO: Rewrite per edit plan Section 6 standard rule (Ch 16).
% Delete: listings. Keep: tables, decision frameworks. Add: one diagram.

% ----------------------------------------------------------
\chapter{Cost Modelling and Budget Control}
\label{ch:cost-modelling}
% TODO: Rewrite per edit plan Section 6 standard rule (Ch 17).
% Delete: listings. Keep: tables, decision frameworks. Add: one diagram.

\partclose{End of Part~3 --- The framework is observable and testable.
Part~4 secures it.}

% ============================================================
% PART 4 — Security and Governance
% (~80 pages, Chapters 18–21)
% ============================================================
\part{Security and Governance}
\label{part:security}

% ----------------------------------------------------------
\chapter{The Agent Threat Landscape}
\label{ch:threat-landscape}
% TODO: Rewrite per edit plan Section 7, Chapter 18.
% Essay-style, no listings.
% Delete: any code fragments that belong in later chapters.
% Add: threat taxonomy table (threat → framework component) + Fig 18.1 (threat-defence mapping Mermaid).

% ----------------------------------------------------------
\chapter{Secure by Design}
\label{ch:secure-by-design}
% TODO: Rewrite per edit plan Section 7, Chapter 19.
% Primary action: listing removals (standard rule).
% Keep: architectural defence patterns.

% ----------------------------------------------------------
\chapter{Enterprise Governance}
\label{ch:enterprise-governance}
% TODO: Rewrite per edit plan Section 7, Chapter 20.
% Primary action: listing removals.
% Add: GDPR right-to-erasure walkthrough using InMemoryAuditLog.allRecordsForUser() and contentHash field.

% ----------------------------------------------------------
\chapter{Runtime Security Monitoring}
\label{ch:runtime-security}
% TODO: Rewrite per edit plan Section 7, Chapter 21.
% Primary action: listing removals.
% Keep: architectural defence patterns.

\partclose{End of Part~4 --- The framework is secure.
Part~5 puts it to work.}

% ============================================================
% PART 5 — Production Agent Architectures
% (~80 pages, Chapters 22–28)
% ============================================================
\part{Production Agent Architectures}
\label{part:production-architectures}

% Each chapter follows the 4-section case study template:
%   §1 Problem and unique constraints
%   §2 Architectural decisions (component diagram)
%   §3 Key implementation points (2–3 excerpts, 15–20 lines each)
%   §4 What to monitor and what to test

% ----------------------------------------------------------
\chapter{The Enterprise Copilot}
\label{ch:enterprise-copilot}
% TODO: Rewrite per edit plan Section 8.
% Add: Fig 22.1 (component composition diagram — AgentRuntime + CoT + TieredMemory + MCP tools + SecurityEnforcer + HITL + OTel).

% ----------------------------------------------------------
\chapter{The Autonomous Research Agent}
\label{ch:research-agent}
% TODO: Rewrite per edit plan Section 8.
% Add: component composition diagram.

% ----------------------------------------------------------
\chapter{The Multi-Agent Pipeline}
\label{ch:multi-agent-pipeline}
% TODO: Rewrite per edit plan Section 8.
% Add: component composition diagram.

% ----------------------------------------------------------
\chapter{The HITL Approval Workflow Agent}
\label{ch:hitl-workflow}
% TODO: Rewrite per edit plan Section 8.
% Add: component composition diagram.

% ----------------------------------------------------------
\chapter{The RAG-Augmented Knowledge Agent}
\label{ch:rag-knowledge-agent}
% TODO: Rewrite per edit plan Section 8.
% Add: component composition diagram.

% ----------------------------------------------------------
\chapter{The Code Generation Agent}
\label{ch:code-generation-agent}
% TODO: Rewrite per edit plan Section 8.
% Add: component composition diagram.

% ----------------------------------------------------------
\chapter{The Autonomous Operations Agent}
\label{ch:operations-agent}
% TODO: Rewrite per edit plan Section 8.
% Add: component composition diagram.

\partclose{End of Part~5 --- Architecture applied.
Part~6 looks ahead.}

% ============================================================
% PART 6 — The Future of Agentic AI
% (~15 pages, Chapters 29–37 — essay style, no code)
% ============================================================
\part{The Future of Agentic AI}
\label{part:future}

% ----------------------------------------------------------
\chapter{Emerging Architectures}
\label{ch:emerging-architectures}
% TODO: Light edit only. Essay-style. No code changes.

% ----------------------------------------------------------
\chapter{The Agent Ecosystem}
\label{ch:agent-ecosystem}
% TODO: Light edit only. Essay-style. No code changes.

% ----------------------------------------------------------
\chapter{Open Problems and Research Directions}
\label{ch:open-problems}
% TODO: Light edit only. Essay-style. No code changes.

% ----------------------------------------------------------
\chapter{Towards Trustworthy Agents}
\label{ch:trustworthy-agents}
% TODO: Light edit only. Essay-style. No code changes.

% ----------------------------------------------------------
\chapter{Agent--Human Collaboration}
\label{ch:agent-human-collaboration}
% TODO: Light edit only. Essay-style. No code changes.

% ----------------------------------------------------------
\chapter{Regulation and Compliance Horizons}
\label{ch:regulation-compliance}
% TODO: Light edit only. Essay-style. No code changes.

% ----------------------------------------------------------
\chapter{The Economics of Agentic Systems}
\label{ch:economics}
% TODO: Light edit only. Essay-style. No code changes.

% ----------------------------------------------------------
\chapter{Closing Perspectives}
\label{ch:closing-perspectives}
% TODO: Light edit only. Essay-style. No code changes.

% ----------------------------------------------------------
\chapter{What Comes Next}
\label{ch:what-comes-next}
% TODO: Light edit only. Essay-style. No code changes.

% ============================================================
\backmatter
% ============================================================

\printindex

\end{document}
